\documentclass[11pt]{article}
\usepackage[margin=1in]{geometry}
\usepackage{amsmath,amssymb,amsthm,hyperref}
\newtheorem{theorem}{Theorem}
\title{A deterministic obstruction to a universal subadditive rate}
\author{ziangni-sys}
\date{Conjecture 00000008492}
\begin{document}
\sloppy
\maketitle
\noindent Both source versions assert that the universal spectrum of rates of
mean convergence for random subadditive sequences is dominated by
$\log n/n$. We disprove this assertion even on a one-point ergodic probability
space. In particular, the obstruction persists when the implied constant and
the starting index may depend arbitrarily on the process.

\begin{theorem}
There is an integrable, nonnegative subadditive process on an ergodic
probability-preserving system whose normalized expectations converge to zero,
but whose errors are not $O(\log n/n)$.
\end{theorem}
\begin{proof}
Let $\Omega=\{*\}$, let $\mu$ be the unit point mass, and let $T$ be the
identity. This system preserves probability and is ergodic: its only subsets
are the empty set and $\Omega$. Define, for every integer $n\geq0$,
\[
 X_n(*)=\sqrt n.
\]
Every $X_n$ is a measurable integrable random variable, is nonnegative, and
$X_0=0$. For all $m,n\geq0$,
\[
 X_{m+n}(*)=\sqrt{m+n}\leq\sqrt m+\sqrt n
 =X_m(*)+X_n(T^m*).
\]
The inequality follows by squaring the nonnegative quantities: the square on
the right is $m+n+2\sqrt{mn}\geq m+n$. Thus this is the usual subadditive
cocycle inequality, with its actual iterate of the measure-preserving map.

Since integration against $\mu$ evaluates at its one point,
\[
 a_n:=\frac{\mathbb E_\mu X_n}{n}
     =\frac{\sqrt n}{n}=\frac1{\sqrt n}\longrightarrow0
 \qquad(n\geq1).
\]
The normalized process has the same limit pointwise and in mean. Its
one-sided error from the limiting value $0$ is exactly $a_n>0$.

For any real constant $C$, the standard logarithm--power limit gives
\[
 C\frac{\log n}{\sqrt n}\longrightarrow0.
\]
Consequently, for every starting index $N$, there exists $n\geq\max(N,1)$
such that $C\log n<\sqrt n$. Dividing by $n>0$ yields
\[
 C\frac{\log n}{n}<\frac{\sqrt n}{n}=a_n.
\]
This excludes every eventual big-$O$ constant, even for this single fixed
process, and proves the theorem.
\end{proof}

\newpage
\section*{Scope of the disproof}
The source concerns random subadditive sequences and does not impose a
uniform bound on all unnormalized $X_n$, mixing assumptions, or a special
symbolic construction. The example is deterministic, which is a valid
special case of a random process, and its underlying probability system is
ergodic. Each random variable is bounded and integrable; the normalized
variables are uniformly bounded by $1$ for $n\geq1$. The unnormalized values
grow as $\sqrt n$. No claim is made about a different hypothesis requiring
$\sup_n\|X_n\|_\infty<\infty$.

The assertion about a universal upper rate is a necessary clause of the
conjecture. Its failure already disproves the conjunction; no assertion about
the separate Sturmian or Thue--Morse extremal constructions is needed.
The counterexample does not merely violate a coefficient-one inequality:
it violates every possible eventual constant multiple of the proposed rate.

\section*{Formal correspondence}
The accompanying Lean project uses Mathlib's actual probability, ergodicity,
integration, filter convergence, and asymptotic definitions.
\begin{itemize}
\item \texttt{mu} (written $\mu$ in the source) is
\texttt{Measure.dirac ()} on \texttt{Unit} with its full sigma-algebra.
An \texttt{IsProbabilityMeasure} instance and \texttt{system\_ergodic}
establish the probability and \texttt{Ergodic T mu} conditions.
\item \texttt{X n} is the actual real square root of the natural index.
\texttt{integrable\_X}, \texttt{X\_nonneg}, \texttt{X\_zero}, and
\texttt{subadditive\_process} verify integrability, nonnegativity,
the zero initial value, and the full iterated subadditive inequality.
\item \texttt{expectation\_X} computes the Bochner integral.
\texttt{mean} divides that integral by the index;
\texttt{mean\_tendsto\_zero} and \texttt{pointwise\_normalized\_limit}
prove convergence, using the real negative-power limit.
\item The logarithm little-$o$ power theorem yields the limit of
$\log n/\sqrt n$. This is formalized as
\texttt{log\_div\_sqrt\_tendsto\_zero}.
The theorem \texttt{no\_eventual\_rate} proves
\[
 \forall C\in\mathbb R\;\forall N\in\mathbb N\;\exists n\geq N,
 \quad n\geq1\quad\hbox{and}\quad
 C\frac{\log n}{n}<a_n.
\]
\item \texttt{not\_bigO} proves the negation of Mathlib's
\texttt{Asymptotics.IsBigO} at \texttt{atTop}, including its norm-based
definition. Positivity identifies the formal error with the one-sided error.
\end{itemize}
Natural-index division at $n=0$ uses Lean's totalized field convention;
all counterexample witnesses have $n\geq1$, so this has no asymptotic effect.

\section*{Reproduction}
The project pins Lean 4.19.0 and Mathlib commit
\begin{center}\small\texttt{c44e0c8ee63ca166450922a373c7409c5d26b00b}.\end{center}
From \texttt{lean/}, run \texttt{lake update}, \texttt{lake exe cache get},
and \texttt{lake build}. Check the complete source with
\begin{center}\small\texttt{lake env lean -DwarningAsError=true Main.lean}.\end{center}
The source prints axiom audits for the principal theorems. Only standard
logical axioms are used; there are no admitted statements or additional axioms.
\end{document}
